\section{Legal \& Trust Module}

\begin{questionbox}[title=Question (a)]
A replica of the exfiltrated object resides physically in \texttt{us-east-1} (United States), while NovaDynamics and its customers are based in the EU. 
\\
To which broader problem of operating in the cloud does this situation correspond? 
\end{questionbox}

\begin{answerbox}
\vspace{3.5cm}
\end{answerbox}

\begin{questionbox}[title=Question (b)]
All of the evidence in Sections 5--6 was provided to you as an export generated by the cloud
provider's own tooling. You have no way to physically verify that this export was not altered
before being handed to you.
\begin{enumerate}[label=(\alph*)]
	\item Why can this reliance be challenged in a court of law, even if the cloud provider
	      acted in good faith?
	\item Based on principles already covered in this course (e.g.\ hashing, chain of custody,
	      repeatability), what would need to be true about how this export was produced for you
	      to defend it as being just as trustworthy as a bit-for-bit disk image acquired
	      directly, as in Laboratory 2? Given what you do \emph{not} control in a cloud
	      environment, is that standard actually achievable here?
\end{enumerate}
\end{questionbox}

\begin{answerbox}
\vspace{4cm}
\end{answerbox}

\begin{questionbox}[title=Question (c)]
Throughout Section 5 you used \texttt{jq}, a general-purpose JSON tool with no forensic
certification, to analyse evidence exports --- and in Section 5.3 you cross-validated one finding
using \texttt{grep} as an independent second method. Reconcile the following: what, precisely,
makes an analysis ``forensically sound'' --- the tool used, or something else? What did the
cross-validation step actually add that simply re-running the same \texttt{jq} query again would
not have?
\end{questionbox}

\begin{answerbox}
\vspace{3cm}
\end{answerbox}